% CHECKPOINT34: isolated source from the complete saved factor-transporter proof.
\statement{A.108}{Proved}{The actual order-three factor transporter in the first central quotient}
Retain A.103--A.107. Let \(e\in\{1,-1\}\) be the integer representing
the unit \(d\bmod3\), and put
\[
 \epsilon=(1-e)/2,\qquad R_e=y^\epsilon r,\qquad
 t=(k+\epsilon)/3\in\widehat{\mathbb Z}.
 \tag{A.108.1}
\]
Then \(R_e^3=1\), its image in the original \(S_3\) quotient is \(r\),
and there exists an actual element \(h\in E_1=F/N_1\) such that
\[
 \boxed{\Psi_1(r)=hR_eh^{-1},\qquad h_{\rm ab}=(t,-t).}
 \tag{A.108.2}
\]
When \(e=-1\), \(R_e=sr^{-1}s\); thus \(hs\), rather than a pure
element, is the corresponding conjugator from \(r^{-1}\).
With the actual matrices
\[
 \begin{gathered}
 J=\begin{pmatrix}0&-1\\1&-1\end{pmatrix},\quad
 D_e=\operatorname{diag}(1,e),\\
 J_e=(-1)^\epsilon Y^\epsilon J=D_eJD_e^{-1},\quad
 H=\begin{pmatrix}1&-b\\0&d^{-1}\end{pmatrix},
 \end{gathered}
 \tag{A.108.3}
\]
the transporter also satisfies
\[
 q(h)J_eq(h)^{-1}=HJH^{-1}.
 \tag{A.108.4}
\]
The signed representation \(q\) is defined on \(E_1\); it is not
being extended to all of \(B_1\). The transporter is not asserted
unique, or to be the image of a specified full-group frame.

\begin{proof}
\emph{A central pro-two kernel and a coprime lifting lemma.}
By A.106 the exact coordinate criterion is
\[
 N_1=N\cap F'\cap\pi_2^{-1}(M_{2,1}).
\]
Consequently
\[
 1\longrightarrow K=(N\cap F')/N_1
 \longrightarrow E_1'\xrightarrow{q}C'\longrightarrow1
 \tag{A.108.5}
\]
is an actual central pro-two extension.
Surjectivity follows from \(q(F')=C'\) for continuous surjections
of compact groups. Centrality follows from the definition of \(N_1\).
The coordinate criterion embeds \(K\) into the pro-two group
\(M_2/M_{2,1}\). No surjectivity of \(N\to M_2\) is needed.

We use the following elementary continuous fact: an action of
\(C_3\) on any pro-two group \(P\) has trivial nonabelian \(H^1\).
For finite \(P\), induct on its order. The nonzero characteristic
elementary abelian subgroup \(\Omega_1(Z(P))\) is central and
action-stable. Kill the projected cocycle by induction; the remaining
cocycle is central. In additive notation it has norm zero, and
\[
 (\rho-1)(\rho+2\rho^2)=3-(1+\rho+\rho^2)
 \tag{A.108.6}
\]
writes it as a coboundary because three is invertible.
For a pro-two group, intersections of the three translates of
open normal subgroups give cofinal action-stable finite quotients.
The compact sets of possible correcting conjugators have the
finite-intersection property, proving the continuous assertion
without a finite-generation hypothesis.
In particular a cocycle in \(E_1'\) whose image is a coboundary in
\(C'\) is itself a coboundary: lift a correcting conjugator in
(A.108.5), then kill the residual \(K\)-cocycle by this lemma.
No \(H^2\) lifting assertion is required.

\emph{The actual cocycle and its abelian normalization.}
The congruences \(d=1+2k\) give \(k\equiv0\pmod3\) if \(e=1\),
and \(k\equiv2\pmod3\) if \(e=-1\); hence \(t\) in (A.108.1)
exists. Multiplication by three on \(\widehat{\mathbb Z}\) is
injective, so this \(t\) is unambiguous.
The marking gives \(yr=sr^{-1}s\), proving all the stated
order and \(S_3\) properties of \(R_e\).

Let \(\rho_e=\operatorname{Ad}(R_e)|_{E_1}\). By A.107,
\[
 \Psi_1(r)=a_eR_e,\qquad
 a_e=u^{-1}y^{-(k+\epsilon)}.
 \tag{A.108.7}
\]
Since \(\Psi_1(r)^3=R_e^3=1\), \(a_e\) is an actual \(C_3\)-cocycle.
On \(F_{\rm ab}=\widehat{\mathbb Z}^{\,2}\), both \(\rho_e\) and
\(\rho=\operatorname{Ad}(r)\) act by
\(\left(\begin{smallmatrix}0&-1\\1&-1\end{smallmatrix}\right)\).
Choose \(h_0=x^t y^{-t}\), viewed in \(E_1\). The equality
\[
 (1-\rho)(t,-t)=(0,-3t)=(0,-k-\epsilon)
\]
therefore makes
\[
 z_e=h_0^{-1}a_e\rho_e(h_0)\in E_1'
 \tag{A.108.8}
\]
a genuine cocycle.

\emph{The signed residual matrix.}
The exact integral identity (A.107.6) is
\[
 q(\rho(g))=Jq(g)J^{-1}(-I)^{\operatorname{exp}_y(g)}.
\]
Both identities defining \(J_e\) in (A.108.3) are literal integral
matrix identities. In particular \((YJ)^3=-I\), whereas \(J_e^3=I\).
On \(E_1'\) the sign character vanishes, so the induced action is
\(q\rho_e=\operatorname{Ad}(J_e)q\).

Put \(Q=q(h_0)\). Since \(\operatorname{exp}_y(h_0)=-t\), the
ordered calculation gives
\[
 q(z_e)J_e=(-1)^{\epsilon-t}Q^{-1}U^{-1}Y^{-k}JQ.
 \tag{A.108.9}
\]
Now \(3t=k+\epsilon\) and \(k\) is even; hence
\(t\equiv\epsilon\pmod2\), and this sign is one.
Using the entire A.103 matrix
\[
 U=\begin{pmatrix}d^{-1}&b\\b&d(1+b^2)\end{pmatrix}
\]
gives the exact equality
\[
 U^{-1}Y^{-k}J
   =\begin{pmatrix}-b&d(b-1-b^2)\\d^{-1}&b-1\end{pmatrix}
   =HJH^{-1}.
 \tag{A.108.10}
\]
Thus, for \(B_e=Q^{-1}HD_e^{-1}\),
\[
 q(z_e)J_e=B_eJ_eB_e^{-1},\qquad
 \det B_e=e\,d^{-1}\equiv1\pmod3.
 \tag{A.108.11}
\]
The sign cancellation in (A.108.9) is needed for this conjugacy.

\emph{Determinant correction at every odd prime.}
Fix an odd prime \(p\). We multiply \(B_{e,p}\) on the right by
a matrix centralizing \(J_e\) whose determinant is
\((\det B_{e,p})^{-1}\).
For \(p\ne3\), the algebra \(\mathbb Z_p[J_e]\) has defining
polynomial \(T^2+T+1\) and unit discriminant \(-3\).
If the polynomial splits modulo \(p\), its distinct roots lift,
giving \(\mathbb Z_p[J_e]\simeq\mathbb Z_p\times\mathbb Z_p\);
the determinant norm is the product and is onto units.
If it is irreducible, the residue field is \(\mathbb F_{p^2}\).
The residue norm onto \(\mathbb F_p^\times\) is onto because the
multiplicative groups are cyclic and the norm exponent is \(p+1\).
Lift a residue unit giving the required determinant modulo \(p\).
The remaining ratio belongs to \(1+p\mathbb Z_p\) and has a scalar
square root: the square-root equation lifts successively since
its derivative is a unit. This scalar correction is also in the
centralizer.
At \(p=3\), the required ratio already belongs to \(1+3\mathbb Z_3\)
by (A.108.11), so a scalar square root alone suffices.
No ramified norm-surjectivity theorem is used.
We obtain \(G_p\in{\rm SL}_2(\mathbb Z_p)\) with
\[
 q_p(z_e)J_e=G_pJ_eG_p^{-1}.
 \tag{A.108.12}
\]

\emph{Conjugators in the actual derived image.}
A.106 gives the actual product
\[
 C=C_{(2)}\times\prod_{p\ {\rm odd}}{\rm SL}_2(\mathbb Z_p).
\]
Its closed derived subgroup is the product of the closed derived
subgroups: it contains each supported derived factor and hence
their dense finite products.
The continuous character computation
\cite[Proposition~5.7(a), printed p.~41]{A_Beyl1986}, together
with character separation, gives
\[
 {\rm SL}_2(\mathbb Z_p)_{\rm ab}=1\ (p\ge5),\qquad
 {\rm SL}_2(\mathbb Z_3)_{\rm ab}=C_3.
 \tag{A.108.13}
\]
Thus \(G_p\) is already in the derived factor for \(p\ge5\).
At three, the abelianization is the original reduction character
\(\chi\) with \(\chi(X)=\chi(Y)=1\). The equality
\(J\bmod3=XYX^{-1}\) gives \(\chi(J)=1\), and \(\chi(-I)=0\).
Therefore
\[
 \chi(J_e)=1+\epsilon=e\quad\text{in }\mathbb F_3.
\]
Right multiplication of \(G_3\) by a suitable power of \(J_e\)
makes its character zero and preserves (A.108.12).
This places every odd-prime conjugator in the actual derived factor.

At two, \(q_2(z_e)\) is a genuine cocycle in the pro-two group
\(C_{(2)}'\) with its actual order-three action.
The coprime lemma gives a correcting conjugator there.
The product of these choices is \(G\in C'\) satisfying
\[
 q(z_e)=G\rho_e(G)^{-1}.
 \tag{A.108.14}
\]
The two-adic factor has not been replaced by all of
\({\rm SL}_2(\mathbb Z_2)\).

Lift \(G\) to \(g\in E_1'\) using (A.108.5). The residual
\(g^{-1}z_e\rho_e(g)\) is a cocycle in the central pro-two group \(K\).
The same coprime lemma gives
\(z_e=j\rho_e(j)^{-1}\) for some \(j\in E_1'\).
Equation (A.108.8) now yields
\[
 a_e=(h_0j)\rho_e(h_0j)^{-1}.
\]
Take \(h=h_0j\). This proves (A.108.2) with its exact abelianization.
Finally \(q(h)=QG\). Equations (A.108.9)--(A.108.10) and
(A.108.14) give (A.108.4).

All arguments take place in the actual compact quotient groups.
Only the existence of a normalized quotient factor transporter
has been proved. Extra quotient centralizers can change its frame;
uniqueness from the full free profinite product is not available here.
\end{proof}

